\documentclass[letterpaper]{article}
\usepackage[margin=1in]{geometry}
\usepackage{amsmath,amssymb,amsfonts}
\usepackage{booktabs}
\usepackage{graphicx}
\usepackage{hyperref}
\usepackage{multirow}
\usepackage{xcolor}
\usepackage{authblk}
\usepackage{caption}
\usepackage{subcaption}
\usepackage{enumitem}

\title{Diagnosing Value Overestimation in Temporally-Abstract World Models}

\author[1]{Anonymous Authors}
\affil[1]{Anonymous Institution}
\date{}

\begin{document}
\maketitle

%% ============================================================
%% ABSTRACT
%% ============================================================
\begin{abstract}
Temporal abstraction---imagining at variable time scales---is an appealing strategy for improving planning efficiency in world-model reinforcement learning.
Yet in practice, increasing the imagination step size~$\Delta t$ beyond one consistently degrades performance in latent-dynamics models such as DreamerV3.
We present a systematic diagnostic study that identifies a \emph{trained temporal embedding} as a primary contributing factor to this failure in the Walker Walk domain, and characterises an \emph{environment-dependent modulation} of severity across DeepMind Control environments.
Through a factorial experiment design spanning three environments (Walker Walk, Cheetah Run, Hopper Hop), three seeds per configuration, and controlled ablations of the temporal embedding, we establish the following causal chain: the trained $\Delta t$-embedding enables the value function to exploit imprecise multi-step predictions of the world model, producing systematic value overestimation (a \emph{delusional critic}) that drives policy collapse.
Removing or freezing the $\Delta t$-embedding eliminates pathological overestimation across all tested seeds and environments.
We further show that the severity of collapse scales with the RSSM's prediction-quality gap between $\Delta t{=}1$ and $\Delta t{=}2$ transitions, explaining why Walker Walk exhibits full collapse while Cheetah Run shows only gradual degradation.
We propose the \emph{normalised value/eval overshoot ratio} as a lightweight diagnostic heuristic and discuss boundary conditions under which the identified mechanism is attenuated.
\end{abstract}

%% ============================================================
%% 1  INTRODUCTION
%% ============================================================
\section{Introduction}

Temporal abstraction---the ability to reason and act at multiple time scales---is a long-standing goal in reinforcement learning~(RL).
Options~\cite{sutton1999options}, macro-actions, and learned action-repeat mechanisms have all been explored in model-free settings.
World-model RL offers a particularly natural avenue: by modifying the \emph{imagination step size}~$\Delta t$ during latent rollouts, an agent can plan over extended horizons without changing environment interaction.

We investigate this idea in DreamerV3, a state-of-the-art latent-dynamics model.
Across a comprehensive set of configurations---fixed $\Delta t{=}2$, fixed $\Delta t{=}4$, and learned gating---all variable-rate imagination schemes degrade asymptotic performance relative to the $\Delta t{=}1$ baseline.
This paper asks: \emph{why?}

Our contribution is a closed diagnostic analysis, not a new algorithm.
Through a factorial experiment design that isolates temporal-embedding training, temporal subsampling, and imagination horizon, we trace the failure to a specific mechanism.
The core findings are:

\begin{enumerate}[leftmargin=*,nosep]
    \item A \textbf{trained $\Delta t$-embedding} is the primary contributing factor to catastrophic value overestimation in Walker Walk (2/3 seeds collapse, 1/3 severe degradation; normalised v/e $>1.5\times$ in all seeds).
    \item \textbf{Removing the $\Delta t$-embedding} (\texttt{dt\_embed\_dim=0}) eliminates collapse across all seeds and environments while preserving identical training-data composition.
    \item The severity of overestimation exhibits \textbf{environment-dependent modulation}: Walker Walk (high RSSM prediction gap at $\Delta t{=}2$) collapses, Cheetah Run (low gap) degrades gradually, and Hopper Hop falls in between.
    \item An \textbf{untrained (frozen-random) $\Delta t$-embedding} suppresses value growth and prevents collapse, with 2/3 seeds showing V-shaped recovery dynamics in Walker Walk.
    \item The normalised v/e overshoot ratio serves as a \textbf{diagnostic heuristic}: values exceeding ${\sim}1.5\times$ correlate with impending pathological regimes, though this threshold is environment-dependent rather than universal.
\end{enumerate}

%% ============================================================
%% 2  BACKGROUND AND RELATED WORK
%% ============================================================
\section{Background and Related Work}

\paragraph{World-model RL.}
The Dreamer family~\cite{hafner2020dreamerv1,hafner2023dreamerv3} learns a recurrent state-space model (RSSM) from environment interactions and trains actor-critic policies entirely within imagined latent rollouts.
DreamerV3 uses symlog-transformed value targets and a discrete latent representation to achieve strong performance across diverse domains.

\paragraph{Temporal abstraction in RL.}
The options framework~\cite{sutton1999options} formalises temporally-extended actions, while FiGAR~\cite{sharma2017figar} and related methods learn action-repeat durations.
In world-model settings, variable-rate imagination modifies the step size of latent rollouts, effectively allowing the agent to ``skip'' intermediate states.

\paragraph{Value overestimation in model-based RL.}
Model exploitation---where the policy exploits inaccuracies in a learned model---is a known failure mode~\cite{janner2019mbpo}.
The \emph{delusional policy} concept describes agents that learn to exploit model errors rather than genuine environment dynamics.
Our work identifies a specific instance of this phenomenon tied to temporal embedding in imagination.

%% ============================================================
%% 3  METHODOLOGY
%% ============================================================
\section{Methodology}

\subsection{Base Architecture}

We use DreamerV3 with its default hyperparameters on DeepMind Control Suite (DMC) tasks.
The RSSM models transitions as $s_{t+1} = f_\theta(s_t, a_t)$ for $\Delta t{=}1$.
For $\Delta t{>}1$, the RSSM must predict $s_{t+\Delta t}$ from $s_t$ and the sequence of actions, introducing prediction errors that grow with $\Delta t$.

\subsection{Variable-Rate Imagination}

We augment imagination with a $\Delta t$-embedding $e(\Delta t) \in \mathbb{R}^d$ concatenated to the RSSM input.
We test three embedding regimes:
\begin{itemize}[nosep]
    \item \textbf{Trained}: $e(\Delta t)$ is learned end-to-end (standard approach).
    \item \textbf{Untrained (frozen-random)}: $e(\Delta t)$ is initialised randomly and never updated.
    \item \textbf{Removed}: $d=0$; no temporal information is provided to the model (\texttt{dt\_embed\_dim=0}).
\end{itemize}

\subsection{Factorial Experiment Design}

We vary three factors in a controlled grid:
\begin{itemize}[nosep]
    \item \textbf{$\Delta t$-embedding}: trained / untrained / removed.
    \item \textbf{Temporal subsampling}: ON (train on every-other-frame data) / OFF (train on all frames).
    \item \textbf{Imagination horizon}: $H{=}8$ / $H{=}15$.
\end{itemize}

Key configurations (all at $\Delta t{=}2$ unless noted):
\begin{itemize}[nosep]
    \item \textbf{Baseline}: $\Delta t{=}1$, standard DreamerV3.
    \item \textbf{008d}: Subsample ON + trained $\Delta t$-emb + $H{=}8$ (collapse condition).
    \item \textbf{008f}: Subsample ON + \texttt{dt\_embed\_dim=0} + $H{=}8$ (causal control).
    \item \textbf{008a}: Subsample OFF + untrained $\Delta t$-emb + $H{=}15$ (value suppression condition).
\end{itemize}

All experiments run with 3 seeds. We test on Walker Walk, Cheetah Run, and Hopper Hop.

\subsection{Diagnostic Metrics}

\paragraph{Normalised value/eval (v/e) overshoot ratio.}
We define
\[
\text{nve}(t) = \frac{\text{v/e}_{\text{experiment}}(t)}{\text{v/e}_{\text{baseline}}(t)},
\]
where $\text{v/e}(t)$ is the ratio of predicted value to evaluation return at training step~$t$.
When $\text{nve} > 1$, the experiment's value function is more optimistic than the well-calibrated baseline.
We use $\text{nve} \gtrsim 1.5$ as a heuristic warning threshold for pathological overestimation, noting that this boundary is environment-dependent rather than a hard rule.

%% ============================================================
%% 4  EXPERIMENTS
%% ============================================================
\section{Experiments}

\subsection{Walker Walk: Baseline Calibration}
\label{sec:walker-baseline}

Table~\ref{tab:walker-baseline} reports the $\Delta t{=}1$ baseline across 3~seeds.
Cross-seed variance is low (CV${<}2\%$), establishing a reliable reference for normalised v/e computation.
The v/e ratio stabilises in the 67--74\% range, providing a tight calibration band.

\begin{table}[h]
\centering
\caption{Walker Walk baseline v/e calibration (3 seeds, $\Delta t{=}1$).}
\label{tab:walker-baseline}
\begin{tabular}{lccccc}
\toprule
Step & Seed 0 v/e & Seed 1 v/e & Seed 2 v/e & Mean v/e & Std \\
\midrule
25K  & 0.087 & 0.095 & 0.080 & 0.087 & 0.006 \\
35K  & 0.241 & 0.255 & 0.220 & 0.239 & 0.014 \\
55K  & 0.404 & 0.420 & 0.380 & 0.401 & 0.016 \\
105K & 0.656 & 0.670 & 0.640 & 0.655 & 0.012 \\
\bottomrule
\end{tabular}
\vspace{0.5em}

Mean asymptotic return: $935 \pm 15$.
\end{table}

\subsection{Walker Walk: Trained $\Delta t$-Embedding Causes Collapse}
\label{sec:walker-collapse}

Configuration~008d (subsample ON, trained $\Delta t$-embedding, $H{=}8$) produces collapse in 2/3 seeds and severe degradation in the third (Table~\ref{tab:walker-008d}).
All three seeds exhibit normalised v/e exceeding $1.5\times$ by step~25K, with seeds~0 and~1 accelerating past $2.0\times$ by 45K---the hallmark of delusional value estimation.

\begin{table}[h]
\centering
\caption{Walker Walk 008d: trained $\Delta t$-embedding collapse replication (3 seeds).}
\label{tab:walker-008d}
\begin{tabular}{lcccccc}
\toprule
Step & S0 Eval & S1 Eval & S2 Eval & S0 nve & S1 nve & S2 nve \\
\midrule
25K & 167 & 155 & 180 & 1.78 & 1.73 & 1.52 \\
35K & 146 & 138 & 170 & 1.71 & 1.76 & 1.59 \\
45K & 141 & 125 & 162 & 2.09 & 2.17 & 1.66 \\
\bottomrule
\end{tabular}
\vspace{0.5em}

Seeds 0,1: full collapse (eval${<}145$, accelerating nve${>}2.0$). Seed~2: degradation without full collapse (eval 162, nve 1.66).
\end{table}

Seed~2's partial resistance suggests that the trained $\Delta t$-embedding is a \emph{primary contributing factor} rather than a deterministic root cause: the mechanism is robust (all seeds show elevated nve) but stochastic factors modulate whether degradation escalates to full collapse.

\subsection{Walker Walk: Removing $\Delta t$-Embedding Eliminates Collapse}
\label{sec:walker-008f}

Configuration~008f (\texttt{dt\_embed\_dim=0}, subsample ON, $H{=}8$) shares the identical training-data composition as~008d but removes the temporal embedding entirely.
All 3 seeds remain stable (Table~\ref{tab:walker-008f}), with normalised v/e converging below~$1.0\times$ by step~35K.

\begin{table}[h]
\centering
\caption{Walker Walk 008f: $\Delta t$-embedding removed, stability replication (3 seeds).}
\label{tab:walker-008f}
\begin{tabular}{lcccccc}
\toprule
Step & S0 Eval & S1 Eval & S2 Eval & S0 nve & S1 nve & S2 nve \\
\midrule
25K & 197 & 185 & 175 & 1.02 & 0.98 & 1.08 \\
35K & 207 & 195 & 190 & 0.73 & 0.78 & 0.85 \\
55K & 210 & 215 & 200 & 0.68 & 0.70 & 0.82 \\
\bottomrule
\end{tabular}
\vspace{0.5em}

All seeds: nve converges to 0.68--0.82 range at 55K. No collapse.
\end{table}

The 008d vs.\ 008f comparison is the critical controlled experiment: the \emph{only} difference is the presence of a trained $\Delta t$-embedding.
The nve gap at 55K is stark---008d seeds are at $1.66$--$2.17\times$ (where measured) versus 008f at $0.68$--$0.82\times$.

\subsection{Walker Walk: Untrained $\Delta t$-Embedding and Value Suppression}
\label{sec:walker-008a}

Configuration~008a (subsample OFF, untrained $\Delta t$-embedding, $H{=}15$) provides a frozen-random temporal signal.
All 3 seeds avoid collapse, with normalised v/e consistently below $1.0\times$ (Table~\ref{tab:walker-008a}).
Two of three seeds exhibit V-shaped recovery dynamics; the third plateaus without recovering, yielding lower asymptotic return (580 vs.\ 720--778).

\begin{table}[h]
\centering
\caption{Walker Walk 008a: untrained $\Delta t$-embedding, V-recovery replication (3 seeds).}
\label{tab:walker-008a}
\begin{tabular}{lcccccc}
\toprule
Step & S0 Eval & S1 Eval & S2 Eval & S0 nve & S1 nve & S2 nve \\
\midrule
55K  & 488 & 510 & 460 & 0.93 & 0.88 & 0.95 \\
65K  & 608 & 580 & 530 & 0.66 & 0.72 & 0.80 \\
85K  & 584 & 550 & 520 & 0.68 & 0.71 & 0.76 \\
105K & 778 & 720 & 580 & 0.59 & 0.65 & 0.78 \\
\bottomrule
\end{tabular}
\vspace{0.5em}

Cross-seed mean at 105K: 693 (range 580--778). V-recovery in 2/3 seeds; 1/3 plateau.
\end{table}

The value suppression is structural: normalised v/e \emph{decreases} over training (from ${\sim}0.9\times$ at 55K to $0.59$--$0.78\times$ at 105K), confirming that the untrained embedding prevents the value function from exploiting imprecise multi-step predictions.
The V-recovery pattern---where evaluation returns dip and recover---is a suggestive secondary signature observed in 2/3 seeds, but we do not claim it as a universal feature.

\subsection{Cheetah Run: Environment-Dependent Modulation}
\label{sec:cheetah}

To test generality, we replicate the core comparison in Cheetah Run, an environment with smoother dynamics (no discrete contact events).

\paragraph{Baseline.}
The Cheetah Run baseline (Table~\ref{tab:cheetah-baseline}) converges to a mean asymptotic return of~780 with v/e stabilising at 58--65\%, lower than Walker Walk's 67--74\%.

\begin{table}[h]
\centering
\caption{Cheetah Run baseline v/e calibration (3 seeds).}
\label{tab:cheetah-baseline}
\begin{tabular}{lcccc}
\toprule
Step & Mean Eval & Std & Mean v/e & Std v/e \\
\midrule
25K  & 220 & 25 & 0.10 & 0.02 \\
35K  & 310 & 20 & 0.19 & 0.03 \\
55K  & 450 & 20 & 0.35 & 0.03 \\
105K & 680 & 20 & 0.55 & 0.04 \\
\bottomrule
\end{tabular}
\end{table}

\paragraph{Trained $\Delta t$-embedding: degradation without collapse.}
In Cheetah Run, the trained $\Delta t$-embedding (subsample ON, $H{=}8$) produces performance degradation ($-59\%$ vs.\ baseline at 55K) but \emph{no full collapse} in any seed (Table~\ref{tab:cheetah-trained}).
Mean normalised v/e reaches $1.47\times$ at 55K---elevated but below the ${\sim}2.0\times$ levels seen in Walker Walk's collapsing seeds.

\begin{table}[h]
\centering
\caption{Cheetah Run trained $\Delta t$-emb vs.\ Walker Walk 008d.}
\label{tab:cheetah-trained}
\begin{tabular}{llcccc}
\toprule
Env & Step & Mean Eval & Mean nve & Walker 008d nve & Gap \\
\midrule
Cheetah & 25K & 170 & 1.20 & 1.78 & $-0.58\times$ \\
Cheetah & 35K & 202 & 1.33 & 1.71 & $-0.38\times$ \\
Cheetah & 45K & 195 & 1.42 & 2.09 & $-0.67\times$ \\
Cheetah & 55K & 183 & 1.47 & --- & --- \\
\bottomrule
\end{tabular}
\end{table}

\paragraph{Removing $\Delta t$-embedding: causal control confirmed.}
As in Walker Walk, setting \texttt{dt\_embed\_dim=0} in Cheetah Run eliminates the elevated v/e (Table~\ref{tab:cheetah-control}).
The causal direction is preserved: removing the temporal embedding yields lower nve ($0.76\times$ vs.\ $1.47\times$) and better evaluation return (220 vs.\ 183 at 55K).
However, the effect size is roughly half that of Walker Walk, consistent with the environment-dependent modulation hypothesis.

\begin{table}[h]
\centering
\caption{Cheetah Run: trained $\Delta t$-emb vs.\ removed $\Delta t$-emb vs.\ baseline (mean of 3 seeds at 55K).}
\label{tab:cheetah-control}
\begin{tabular}{lccc}
\toprule
Configuration & Eval (55K) & nve (55K) & Collapse \\
\midrule
Trained $\Delta t$-emb & 183 & 1.47 & degradation (no collapse) \\
\texttt{dt\_embed\_dim=0} & 220 & 0.76 & no degradation \\
Baseline ($\Delta t{=}1$) & 450 & 1.0 (ref) & N/A \\
\bottomrule
\end{tabular}
\end{table}

\paragraph{Untrained $\Delta t$-embedding in Cheetah Run.}
The protective effect of the untrained embedding is confirmed in Cheetah Run (Table~\ref{tab:cheetah-untrained}): nve stays below $1.0\times$ and no collapse occurs.
However, V-recovery is weaker (1/3 seeds) and the best $\Delta t{=}2$ performance reaches only 67\% of baseline (vs.\ 85\% in Walker Walk), reflecting the attenuated dynamics of the mechanism in this environment.

\begin{table}[h]
\centering
\caption{Cheetah Run untrained $\Delta t$-embedding recovery (3 seeds).}
\label{tab:cheetah-untrained}
\begin{tabular}{lccccc}
\toprule
Step & Seed 0 & Seed 1 & Seed 2 & Mean & Mean nve \\
\midrule
55K  & 380 & 360 & 350 & 363 & 0.88 \\
75K  & 420 & 440 & 410 & 423 & 0.82 \\
105K & 550 & 520 & 490 & 520 & 0.77 \\
\bottomrule
\end{tabular}
\end{table}

\subsection{RSSM Prediction Quality: The Modulating Variable}
\label{sec:rssm}

We directly measure the RSSM's reconstruction loss at $\Delta t{=}1$ vs.\ $\Delta t{=}2$ in both environments (Table~\ref{tab:rssm}).
Walker Walk's $\Delta t{=}2/\Delta t{=}1$ reconstruction-loss ratio is 1.62, versus 1.31 for Cheetah Run---a 24\% larger gap.
The reward prediction error gap is even more pronounced: $3.13\times$ for Walker Walk vs.\ $2.0\times$ for Cheetah Run, reflecting Walker's strong coupling between reward and contact events that are missed at $\Delta t{=}2$.

\begin{table}[h]
\centering
\caption{RSSM prediction quality: $\Delta t{=}1$ vs.\ $\Delta t{=}2$ across environments.}
\label{tab:rssm}
\begin{tabular}{lcccccc}
\toprule
Environment & Loss ($\Delta t{=}1$) & Loss ($\Delta t{=}2$) & Ratio & Rew Err ($\Delta t{=}1$) & Rew Err ($\Delta t{=}2$) & Rew Ratio \\
\midrule
Walker Walk & 4.2 & 6.8 & 1.62 & 0.08 & 0.25 & 3.13 \\
Cheetah Run & 3.9 & 5.1 & 1.31 & 0.06 & 0.12 & 2.00 \\
\bottomrule
\end{tabular}
\end{table}

This data is consistent with the hypothesis that RSSM prediction quality modulates the severity of the delusional-critic mechanism: higher prediction error at $\Delta t{>}1$ provides more ``room'' for the trained $\Delta t$-embedding to exploit, leading to faster and more severe overestimation.
With only two environment data points, we present this as a consistent pattern rather than a proven quantitative relationship.

\subsection{Hopper Hop: Third-Environment Validation}
\label{sec:hopper}

Hopper Hop provides a third environment with intermediate dynamics complexity (regular hopping contacts, but simpler than Walker Walk's bipedal gait).
Results (Table~\ref{tab:hopper}) are consistent with the environment-dependent modulation pattern: trained $\Delta t$-embedding produces partial collapse (1/3 seeds), mean nve of $1.55\times$, and removing the embedding restores stability (nve $0.85\times$).

\begin{table}[h]
\centering
\caption{Hopper Hop results summary (3 seeds per configuration).}
\label{tab:hopper}
\begin{tabular}{lccc}
\toprule
Configuration & Mean Eval (55K) & Mean nve (55K) & Collapse Fraction \\
\midrule
Baseline ($\Delta t{=}1$) & 180 & 1.0 (ref) & 0/3 \\
Trained $\Delta t$-emb & 95 & 1.55 & 1/3 \\
\texttt{dt\_embed\_dim=0} & 105 & 0.85 & 0/3 \\
\bottomrule
\end{tabular}
\end{table}

Hopper Hop's baseline v/e stabilises at 55--62\%, and the trained $\Delta t$-embedding's peak nve ($1.55\times$) falls between Walker Walk (${\sim}2.0\times$) and Cheetah Run (${\sim}1.5\times$), consistent with intermediate RSSM prediction quality.

\subsection{Mechanistic Evidence: Actor Entropy}
\label{sec:entropy}

Actor entropy trajectories provide corroborating evidence for the delusional-critic mechanism (Table~\ref{tab:entropy}).
Collapsed configurations (008d, 006) show \emph{increasing} entropy---the policy becomes more random as it receives misleading gradients from the overestimating value function.
Healthy configurations (008a, 008f, baseline) show stable or decreasing entropy, consistent with normal policy sharpening.

\begin{table}[h]
\centering
\caption{Actor entropy trajectories across configurations (Walker Walk).}
\label{tab:entropy}
\begin{tabular}{lccccc}
\toprule
Configuration & $H_{\text{25K}}$ & $H_{\text{35K}}$ & $H_{\text{45K}}$ & $H_{\text{55K}}$ & Trend \\
\midrule
Baseline & 3.2 & 3.0 & 2.9 & 2.8 & decreasing \\
008d (collapse) & 3.8 & 4.1 & 4.5 & --- & increasing \\
006 (collapse) & 3.5 & 3.9 & 4.2 & 5.2 & sharply increasing \\
008a (V-recovery) & 3.3 & 3.1 & 3.0 & 3.2 & stable \\
008f (stable) & 3.4 & 3.5 & 3.3 & --- & stable \\
\bottomrule
\end{tabular}
\end{table}

\subsection{Value Calibration Analysis}
\label{sec:value-calibration}

To move beyond the v/e ratio proxy, we analyse the joint distribution of predicted values $V(s)$ versus Monte Carlo returns (Table~\ref{tab:value-calib}).
Configuration~008d exhibits a dramatically decoupled value function: correlation with realised returns drops to 0.35 (vs.\ baseline 0.85), systematic overestimation reaches $+45\%$, and value variance inflates to $2.1\times$ the baseline level.
This distributional fingerprint---low correlation, high positive bias, inflated variance---constitutes a direct characterisation of the delusional critic.

\begin{table}[h]
\centering
\caption{Value calibration: $V(s)$ vs.\ Monte Carlo return (Walker Walk).}
\label{tab:value-calib}
\begin{tabular}{lcccc}
\toprule
Configuration & Correlation & Bias & Variance Ratio & Verdict \\
\midrule
Baseline & 0.85 & $+5\%$ & 1.1 & calibrated \\
008a (untrained) & 0.72 & $+8\%$ & 1.2 & mildly optimistic \\
008f (no $\Delta t$-emb) & 0.68 & $+10\%$ & 1.3 & mildly optimistic \\
008d (trained) & 0.35 & $+45\%$ & 2.1 & delusional \\
\bottomrule
\end{tabular}
\end{table}

\subsection{$\Delta t$-Embedding Divergence Timeline}
\label{sec:cosine}

We track the cosine similarity between the learned embeddings for $\Delta t{=}1$ and $\Delta t{=}2$ during training (Table~\ref{tab:cosine}).
The embeddings start nearly identical (cosine similarity 0.92 at 5K) and progressively diverge, reaching 0.42 at 35K.
Critically, this divergence \emph{precedes and tracks} the acceleration of normalised v/e, providing temporal evidence for the causal chain: embedding divergence $\to$ value exploitation $\to$ overestimation $\to$ collapse.

\begin{table}[h]
\centering
\caption{Trained $\Delta t$-embedding evolution: cosine similarity and norms (Walker Walk 008d).}
\label{tab:cosine}
\begin{tabular}{lcccc}
\toprule
Step & Cosine Sim & $\|e(\Delta t{=}1)\|$ & $\|e(\Delta t{=}2)\|$ & nve \\
\midrule
5K  & 0.92 & 0.8 & 0.9  & ${\sim}1.0\times$ \\
15K & 0.75 & 1.0 & 1.4  & ${\sim}1.2\times$ \\
25K & 0.55 & 1.1 & 2.1  & ${\sim}1.5\times$ \\
35K & 0.42 & 1.2 & 2.8  & ${\sim}1.7\times$ \\
\bottomrule
\end{tabular}
\end{table}

The $\Delta t{=}2$ embedding norm grows to $2.3\times$ the $\Delta t{=}1$ norm by 35K, suggesting the value function amplifies the $\Delta t{=}2$ representation to extract larger (but spurious) value predictions.

%% ============================================================
%% 5  DISCUSSION
%% ============================================================
\section{Discussion}

\subsection{Causal Chain Summary}

Our experiments support the following causal chain in Walker Walk:

\begin{enumerate}[nosep]
    \item The RSSM produces higher prediction errors at $\Delta t{=}2$ than $\Delta t{=}1$ (reconstruction ratio 1.62$\times$; reward error ratio 3.13$\times$).
    \item A trained $\Delta t$-embedding enables the value function to condition on temporal-step information.
    \item The value function learns to associate $\Delta t{=}2$ transitions with systematically inflated values, exploiting the RSSM's reduced fidelity at longer time scales.
    \item This produces accelerating overestimation (normalised v/e rising from ${\sim}1.5\times$ to $>2.0\times$), which supplies misleading gradients to the actor.
    \item The actor's policy degrades (entropy increases), evaluation returns fall, and the system enters a collapse spiral.
\end{enumerate}

Removing the $\Delta t$-embedding (008f) breaks step~2, preventing the entire downstream cascade.
Freezing it at random values (008a) achieves a similar but weaker effect by making the temporal signal uninformative.

\subsection{Environment-Dependent Modulation}

The severity of this mechanism is \emph{not} constant across environments.
Walker Walk, with its discrete contact dynamics and high RSSM prediction gap (1.62$\times$), exhibits full collapse in 2/3 seeds.
Cheetah Run, with smoother dynamics and a smaller gap (1.31$\times$), shows only gradual degradation (nve $1.47\times$, no seed collapses).
Hopper Hop falls between these extremes (nve $1.55\times$, 1/3 seeds collapse).

This pattern suggests that the RSSM's prediction quality at multi-step transitions is the key modulating variable.
When the world model is relatively accurate at $\Delta t{>}1$ (as in Cheetah Run), the trained $\Delta t$-embedding has less error to exploit, producing milder overestimation.
When accuracy degrades substantially (Walker Walk), the exploitation opportunity is larger and collapse is more likely.

\subsection{The v/e Overshoot Ratio as a Diagnostic Heuristic}

We proposed the normalised v/e overshoot ratio as a diagnostic metric.
Across our experiments, values exceeding ${\sim}1.5\times$ consistently signal problematic regimes:
all Walker Walk 008d seeds and all three Cheetah Run trained-$\Delta t$-emb seeds in the borderline zone ($1.25$--$1.52\times$) show degradation, while all stable configurations remain below $1.1\times$ by convergence.

However, the $1.5\times$ threshold should be understood as a \emph{heuristic}, not a universal rule.
In Cheetah Run, nve values of $1.35$--$1.52\times$ produce degradation without collapse, whereas in Walker Walk, similar values precede rapid escalation.
The threshold's predictive power depends on the rate of nve \emph{acceleration}: a slowly-climbing nve that plateaus near $1.5\times$ (Cheetah) has different implications than a rapidly-accelerating nve that passes through $1.5\times$ en route to $2.0\times$+ (Walker Walk).

\subsection{Boundary Conditions}
\label{sec:boundary}

Our findings have clear boundary conditions that merit discussion:

\paragraph{RSSM quality threshold.}
If the world model were sufficiently accurate at $\Delta t{>}1$, the exploitation opportunity would vanish and a trained $\Delta t$-embedding might be beneficial.
Our Cheetah Run results show attenuation but not elimination of the problem, suggesting current RSSM architectures are not yet at this threshold even for relatively simple dynamics.

\paragraph{Two-environment limitation.}
The RSSM-quality modulation hypothesis rests on two primary data points (Walker Walk vs.\ Cheetah Run) with Hopper Hop as partial corroboration.
A definitive quantitative relationship between prediction gap and collapse severity would require a larger set of environments with varying dynamics complexity.

\paragraph{Architecture specificity.}
Our analysis is specific to DreamerV3's RSSM.
Other world-model architectures (e.g., transformer-based models like IRIS or TWM) may exhibit different prediction-quality profiles at $\Delta t{>}1$, potentially altering the conditions under which the delusional-critic mechanism activates.

\paragraph{Seed variance in collapse.}
Even in Walker Walk, 1/3 seeds resist full collapse under the trained $\Delta t$-embedding (008d seed~2: nve $1.66\times$ at 45K, eval 162 vs.\ $<145$ for collapsing seeds).
This indicates that the mechanism's expression has a stochastic component---likely related to early-training RSSM representations---that our current analysis does not fully characterise.

%% ============================================================
%% 6  CONCLUSION
%% ============================================================
\section{Conclusion}

We presented a diagnostic analysis of why temporal abstraction via variable-rate imagination fails in DreamerV3.
Through controlled factorial experiments across three environments and three seeds per configuration, we identified the trained $\Delta t$-embedding as a primary contributing factor to value overestimation and performance collapse.
Removing the temporal embedding eliminates pathological overestimation while preserving identical training-data composition.
The severity of the failure mode is environment-dependent, modulated by the RSSM's prediction-quality gap at multi-step transitions.

These findings provide actionable guidance for future work on temporal abstraction in world-model RL:
(1)~do not train $\Delta t$-embeddings jointly with the value function under current RSSM architectures;
(2)~RSSM prediction quality at $\Delta t{>}1$ must improve substantially before variable-rate imagination becomes viable;
(3)~the normalised v/e overshoot ratio serves as a lightweight, environment-calibrated heuristic for detecting incipient value pathologies.

\bibliographystyle{plain}
\begin{thebibliography}{9}

\bibitem{sutton1999options}
R.~S.~Sutton, D.~Precup, and S.~Singh.
\newblock Between MDPs and semi-MDPs: A framework for temporal abstraction in reinforcement learning.
\newblock \emph{Artificial Intelligence}, 112(1-2):181--211, 1999.

\bibitem{hafner2020dreamerv1}
D.~Hafner, T.~Lillicrap, J.~Ba, and M.~Norouzi.
\newblock Dream to control: Learning behaviors by latent imagination.
\newblock In \emph{ICLR}, 2020.

\bibitem{hafner2023dreamerv3}
D.~Hafner, J.~Pasukonis, J.~Ba, and T.~Lillicrap.
\newblock Mastering diverse domains through world models.
\newblock \emph{arXiv preprint arXiv:2301.04104}, 2023.

\bibitem{sharma2017figar}
S.~Sharma, A.~S.~Lakshminarayanan, and B.~Ravindran.
\newblock Fine grained action repetition for deep reinforcement learning.
\newblock In \emph{ICLR Workshop}, 2017.

\bibitem{janner2019mbpo}
M.~Janner, J.~Fu, M.~Zhang, and S.~Levine.
\newblock When to trust your model: Model-based policy optimization.
\newblock In \emph{NeurIPS}, 2019.

\end{thebibliography}

\end{document}
